% =====================================================================================================
% Chương 4, PHẦN B (mục 4.6–4.11) — bản viết lại, phiên bản 3 (2026-10-12). ĐÂY LÀ BẢN NỘP.
% Bản 1 (09-chuong4-phanB.tex, 4.778 tiếng văn xuôi) và bản 2 (09-chuong4-phanB-v2.tex, 3.949 tiếng) được giữ
% nguyên; cả hai vượt ngân sách trang. Bản 3 giữ cùng số liệu, bảng, hình và nhãn, chỉ rút lời văn và tách câu dài
% (09-chuong4-phanB-NOTES.md, mục 1 và 7).
% Nối ngay sau 08-chuong4-phanA.tex trong cùng \input{noi-dung/08-chuong4}; không có \chapter mới.
% Mỗi dòng "% EV:" trỏ tới EVIDENCE-REGISTER.md (A01…K52), EVIDENCE-REGISTER-ADDENDUM-1.md (N-01…N-30,
% L1a-01…09, mục 0.x), EVIDENCE-REGISTER-ADDENDUM-2-CH3.md (X3-xx) hoặc EVIDENCE-REGISTER-ADDENDUM-3-CH4B.md
% (N4B-01…N4B-14, tạo cùng phần này).
% Không dùng gói mới: hình vẽ bằng TikZ thường; danh sách dùng enumitem như chương 3.
% =====================================================================================================

% -----------------------------------------------------------------------------------------------------
\section{Vai trò của trần và thứ tự ưu tiên chọn phương án}
\label{muc:c4-vaitro}

Điều phối không trần và luật C trần 30 giây dùng cùng tệp cấu hình và cùng bộ chọn ưu tiên ổn định, chỉ khác nhau
ở trần giờ trả. Chênh lệch giữa chúng ở các mục trước vì vậy là chênh lệch khi thêm trần mà giữ nguyên bộ chọn, trên
hai tập khách khác nhau. Câu hỏi còn mở là đánh đổi này có còn khi bộ chọn ưu tiên chi phí thay vì ổn định (tiêu chí
4 thay cho tiêu chí 3 ở Bảng~\ref{bang:rangbuoc}) hay không.
% EV: A08; H01; NEW:12 (tệp U và C-30 chỉ khác /policies/0/limits/1/hardLimit); C04; A14; X3-07; H02

Bằng chứng hiện có là một phép sàng lọc trên tập tắc đường cố định. Ở mỗi lượt quyết định đã ghi, bộ chọn được chạy
lại trên các phương án đã qua phép kiểm của cấu hình đó, sau khi bỏ tiêu chí 3; mô phỏng không chạy tiếp từ lựa chọn
mới. Với thứ tự gốc, phép chạy lại cho đúng lựa chọn đã ghi ở cả 62.837 lượt quyết định.
% EV: N4B-01 (r17_wp14r_check.py:1-19, 32; WP17:3-22); H03 (PV:2-6; tổng của năm cấu hình); H10

\begin{table}[htbp]
\centering
\caption{Phép sàng lọc bỏ tiêu chí ổn định của bộ chọn trên tập tắc đường cố định (16 kịch bản mỗi cấu hình, bộ
chạy cũ). Mỗi cột dùng các lượt quyết định đã ghi của chính cấu hình đó, nên hai cột là hai tập trạng thái khác
nhau. Chi phí kế hoạch là thời lượng còn lại của các tuyến dự kiến, không phải chi phí vận hành.}
\label{bang:c4-sangloc}
\small
\begin{tabular}{@{}p{10.6cm}rr@{}}
\toprule
\textbf{Chỉ số} & \textbf{U} & \textbf{C-30}\\
\midrule
Lượt quyết định có yêu cầu mới & 1.714 & 1.715\\
Lượt đổi lựa chọn khi bỏ tiêu chí ổn định & 445 (25,96\%) & 104 (6,06\%)\\
% EV: H04 (445/1.714; 104/1.715; số yêu cầu được nhận không đổi)
\midrule
\multicolumn{3}{@{}l}{\textit{Chi phí kế hoạch giảm (\%), ba cách tính}}\\
Trung bình trên các lượt đổi lựa chọn & 7,79 & 8,00\\
Trung bình trên mọi lượt có yêu cầu mới & 2,02 & 0,49\\
Tổng phần giảm chia tổng chi phí kế hoạch đã chọn, trên mọi lượt quyết định & 0,3069 & 0,0794\\
% EV: H05 (NEW:6; PV:2,5); K15; A16; I09
\midrule
\multicolumn{3}{@{}l}{\textit{Tổng dời giờ trả do điều phối tăng thêm trên đội xe, ở các lượt đổi lựa chọn (s)}}\\
Trung vị & 90,8 & 12,5\\
Lớn nhất & 390,8 & 35,5\\
% EV: H06 (WP17:8,22; PV:2,5); K17
\bottomrule
\end{tabular}
\end{table}

Cách đọc bảng: với điều phối không trần, bỏ tiêu chí ổn định làm lựa chọn đổi ở 445 trong 1.714 lượt có yêu cầu
mới, và ở các lượt đó chi phí kế hoạch giảm trung bình 7,79\%.
% EV: H04; H05

Số yêu cầu được nhận không đổi ở lượt nào. Tỉ lệ lượt đổi lựa chọn là 25,96\% khi không trần, 18,95\% với trần 120
giây, 10,08\% với trần 60 giây và 6,06\% với trần 30 giây. Mỗi tỉ lệ đo trên trạng thái của chính cấu hình, nên đây
không phải tác động của trần trên cùng một trạng thái; kết quả chỉ gợi ý rằng dưới trần 30 giây, bộ chọn có ít cơ
hội đổi hơn. Chi phí kế hoạch tính lặp qua các lượt, nên ba con số chi phí không cộng thành mức tiết kiệm vận hành.
Dưới trần 30 giây, lượt có tổng dời tăng thêm lớn nhất, 35,5 giây, chia cho hai khách, mỗi khách 17,77 giây.
% EV: H04 (325/1.715; 173/1.716); K16; H05; K15; A16; I09; H06; K17 (CASE:1-18: hai khách, mỗi khách 17.770 ms)

Phép thử đổi thứ tự ưu tiên của bộ chọn trên 34 kịch bản chính chưa được chạy; bằng chứng hiện có là phép sàng lọc
trên trạng thái cố định ở trên, vốn không cho biết kết quả phục vụ và đúng giờ trên quỹ đạo. Phép sàng lọc cũng không
chuyển sang tập kiểm tra chính, vì nó dùng thời gian đi không đổi và bộ chạy cũ. Khóa luận vì vậy chưa tách được phần
của trần và phần của bộ chọn trong độ tin cậy đã đo; chương kết luận đặt phép thử này lên đầu các hướng phát triển.
% EV: H12 (chưa chạy); REG mục L (L02–L05 chưa có kết quả); H10; NEW:13; B17; I12; H02; K18

% -----------------------------------------------------------------------------------------------------
\section{Lời hứa theo dự báo}
\label{muc:c4-dubao}

Trên 21 kịch bản ngày thường có dự báo, số khách của luật C trần 30 giây tới muộn hơn lời hứa đầu quá một phút
giảm từ 68 xuống 2, trong 1.471 khách được phục vụ (gộp theo khách). Tính trung bình theo kịch bản, tỉ lệ này giảm
từ 6,01\% xuống 0,17\%.
% EV: F04 (6,01% → 0,17% trung bình theo ô; gộp 2/1.471); F17 (2 so với 68); RP1:39-40; K24

Lời hứa theo dự báo được tính lại trên các lần chạy đã ghi (mục~\ref{muc:dubao}); trễ được so với lời hứa đầu của
chính cách hứa. Đối chứng thứ hai là \emph{đệm cùng giá}: mỗi lời hứa đầu được kéo dài một tỉ lệ cố định của thời
gian còn lại tới giờ trả, chọn theo từng khung giờ để tổng phần kéo dài bằng của dự báo (Bảng~\ref{bang:c4-dubao}).
% EV: F03 (chỉ vế phát lại); F01; N4B-03 (forecast_family_a.py:15-17: lời hứa đệm = P0 + k_w(P0 − t0); FA:2)

\begin{table}[htbp]
\centering
\caption{Khách của luật C trần 30 giây tới muộn hơn lời hứa đầu quá 60 giây theo ba cách hứa, trên 21 kịch bản
ngày thường có dự báo (1.471 khách được phục vụ). Mọi số là trung bình theo kịch bản, tính bằng \%; dòng cuối là
tỉ lệ khách tới sớm hơn lời hứa quá 60 giây.}
\label{bang:c4-dubao}
\small
\begin{tabular}{@{}lrrr@{}}
\toprule
\textbf{Nhóm kịch bản} & \textbf{Không dự báo} & \textbf{Đệm cùng giá} & \textbf{Dự báo}\\
\midrule
Cả 21 kịch bản & 6,01 & 1,94 & 0,17\\
Khung 7 giờ (8 kịch bản) & 14,44 & 4,24 & 0,27\\
Khung 8 giờ (8 kịch bản) & 1,35 & 0,86 & 0,18\\
Khung 17 giờ (5 kịch bản) & 0 & 0 & 0\\
% EV: F04; N4B-03 (FR:40-45; FA:8-11: cột đệm 1,94 / 4,24 / 0,86 / 0)
\midrule
Tới sớm quá 60 s, cả 21 kịch bản & 0 & chưa đo & 0,08\\
% EV: N4B-03 (FR:46; FA:8: E naive 0, disp 0,0008; phép đệm không tính chỉ số này)
\bottomrule
\end{tabular}
\end{table}

Cách đọc bảng: ở khung 7 giờ, trung bình mỗi kịch bản có 14,44\% khách tới muộn hơn lời hứa đầu quá một phút khi
không dự báo, 4,24\% khi đệm cùng giá và 0,27\% khi có dự báo.
% EV: F04; N4B-03

Khoảng tin cậy 95\% tính theo phân phối $t$ trên 8 cụm ngày và khung giờ cho chênh lệch dự báo trừ không dự báo là
từ $-10{,}54$ đến $-1{,}15$ điểm phần trăm. Với dự báo trừ đệm cùng giá, chênh lệch là $-1{,}77$ điểm phần trăm và
khoảng chạy từ $-3{,}62$ đến $+0{,}07$, tức chạm 0. Lợi thế so với cách đệm này vì vậy còn mong manh, và chỉ dựa trên
ba buổi sáng.
% EV: F05 (FA:33,36; FR:51-60); I05

Lợi ích này không riêng của luật C. Luật hạn chót và luật trần độ dời khách thấy, cùng ở mức 30 giây, cũng xuống
0,17\%; điều phối không trần giảm từ 9,72\% xuống 3,90\% (trung bình theo kịch bản). Ở khung 17 giờ, tỉ lệ 6,76\% của
điều phối không trần không đổi khi có dự báo, còn luật C ở khung này đã là 0\% từ trước. Dự báo và trần nhắm vào hai
phần khác nhau của độ lệch giờ trả.
% EV: F08 (FR:65,79-86; FA:3,6,13,23); K37; FA:11 (C-30 khung 17 giờ 0)

Giá của dự báo là lời hứa dài hơn: trên 21 kịch bản, lời hứa trả dài thêm trung bình 13,7 giây (khung 7 giờ 27,8
giây) và lời hứa đón 3,2 giây, trung bình theo kịch bản. Con số 7,37 giây ở Bảng~\ref{bang:c4-tran120} nhỏ hơn vì nó
chia cho cả 2.338 khách, kể cả khách cuối tuần có phần dài thêm bằng 0. Vì phép kiểm khả thi vẫn dùng giờ không dự
báo, 34 trong 1.471 khách có giờ đón hiển thị muộn hơn giờ đón muộn nhất (mục~\ref{muc:dubao}).
% EV: F06 (FR:67-69; FA:8-9); C09; N-16; N4B-10 (trung bình gộp gồm khách cuối tuần, phần dài thêm 0); F02; N-28; F07

Tắc đường trong mô phỏng là một hệ số chung cho cả thành phố, nên kéo dãn lời hứa theo hệ số thật cho giờ trả khớp
trong vòng 2 giây với mọi khách không bị điều phối dời (2.305 trên 2.305 khách ở 34 kịch bản). Vì vậy 0,17\% là mức
một dự báo theo giờ trong ngày đạt được trong thế giới gần như tất định này, không phải hiệu quả ngoài đời.
% EV: F10 (FR:77; RX:8; FA:28); I02; B06

Hình~\ref{hinh:c4-so} chia số giây trễ của 68 khách nói trên theo sổ (công thức~\eqref{eq:so}): phần lớn được hồ sơ
lịch sử giải thích, còn ở chiều tuyết (mục~\ref{muc:c4-tuyet}) thì gần như ngược lại. Phần thứ hai phụ thuộc dự báo
đang dùng, và phần còn lại không phải phần ``không thể biết trước''.
% EV: F09 (FR:72-75; FA:30; RX:3); N4B-05 (chiều tuyết); K04; K19

\begin{figure}[htbp]
\centering
\begin{tikzpicture}[x=1cm,y=1cm,font=\footnotesize,
  nhan/.style={font=\scriptsize,inner sep=1.5pt}]
% Thanh trên: ngày thường. Độ dài 11 cm = 6.217,4 s; ranh giới 0,1445 và 9,6145 cm (FA:30).
\draw[fill=black!75] (0,1.6) rectangle (0.1445,2.2);
\draw[fill=black!20] (0.1445,1.6) rectangle (9.6145,2.2);
\draw[fill=white] (9.6145,1.6) rectangle (11,2.2);
\node at (4.88,1.9) {5.353 s (86,1\%)};
\node[align=center,nhan] at (10.307,1.9) {783 s\\(12,6\%)};
\draw (0.072,2.2)--(0.072,2.45);
\node[nhan,anchor=south west] at (0.0,2.45) {82 s (1,3\%)};
\node[anchor=east,align=right] at (-0.15,1.9) {Ngày thường, 21 kịch bản\\68 khách, 6.217 s};
% Thanh dưới: chiều tuyết. Độ dài 11 cm = 3.445,9 s; ranh giới 0,0273 và 0,8302 cm (NEW:5).
\draw[fill=black!75] (0,0) rectangle (0.0273,0.6);
\draw[fill=black!20] (0.0273,0) rectangle (0.8302,0.6);
\draw[fill=white] (0.8302,0) rectangle (11,0.6);
\node at (5.915,0.3) {3.185,9 s (92,5\%)};
\draw (0.014,0)--(0.014,-0.2);
\node[nhan,anchor=north west] at (-0.02,-0.2) {8,6 s};
\draw (0.43,0)--(0.43,-0.62);
\node[nhan,anchor=north west] at (0.39,-0.62) {251,5 s (7,3\%)};
\node[anchor=east,align=right] at (-0.15,0.3) {Chiều tuyết, 8 kịch bản\\41 khách, 3.445,9 s};
% Chú giải
\draw[fill=black!75] (0,-1.55) rectangle (0.3,-1.3);
\node[anchor=west] at (0.35,-1.425) {phần tính cho điều phối};
\draw[fill=black!20] (3.7,-1.55) rectangle (4.0,-1.3);
\node[anchor=west] at (4.05,-1.425) {phần được hồ sơ lịch sử giải thích};
\draw[fill=white] (8.6,-1.55) rectangle (8.9,-1.3);
\node[anchor=west] at (8.95,-1.425) {phần còn lại};
\end{tikzpicture}
\caption{Ba phần của số giây trễ theo sổ theo từng khách (công thức~\eqref{eq:so}), với luật C trần 30 giây và các
khách tới muộn hơn lời hứa đầu không dự báo quá 60 giây: 68 khách của 21 kịch bản ngày thường và 41 khách của 8
kịch bản chiều tuyết 15/11/2018. Mỗi thanh dài bằng 100\% số giây trễ của nhóm đó. Phần được hồ sơ lịch sử giải
thích phụ thuộc dự báo đang dùng; phần còn lại là phần dự báo này chưa giải thích được, không phải phần không thể
biết trước.}
\label{hinh:c4-so}
\end{figure}
% EV (ngày thường): F09 (FA:30: 81,7 + 5.352,6 + 783,0 = 6.217,4 s; RX:3: 1,3 / 86,1 / 12,6%)
% EV (chiều tuyết): F13 (8,6 trên 3.445,9 s); N4B-05 (NEW:5: 8,566 / 251,496 / 3.445,933 s; RP1:24-26: 7,3%; SNOW:10: 92,5%; FR:104,186)
% Tọa độ: x = giây / tổng × 11 cm (tính tay; xem NOTES mục 4).

\textbf{Đưa dự báo vào bộ điều phối.} Cách này được thử với luật C và luật hạn chót, cùng ở mức 30 giây, trên 16
kịch bản khung 7 và 8 giờ. Có 3 trong 32 lần chạy dừng, vì xe trong FleetPy tới một điểm đón muộn hơn giờ đón muộn
nhất 1,67 hoặc 2,69 giây. Ở các lần chạy xong, luật C không phải phục hồi lần nào, còn luật hạn chót vẫn phải phục
hồi ở 47,6\% lượt quyết định, so với 62,8\% khi không đưa dự báo vào (trung bình trên 15 kịch bản). Tỉ lệ khách bị
luật hạn chót phá lời hứa ở khung 7 giờ giảm từ 39,0\% xuống 11,0\%, nên con số ở Bảng~\ref{bang:c4-nghiavu} phụ
thuộc cách hứa. Vì có lần chạy dừng, cách triển khai được chọn là lớp lời hứa.
% EV: F12 (FR:117-146; RX:9-15); X3-29; N4B-07 (FR:29,34,120-121; RX:9,11,13-15); C14 (lưu ý: phụ thuộc cách hứa)

\textbf{Học máy.} Ba mô hình cố định ước lượng độ trễ của từng khách từ các đặc trưng có lúc hứa: cây tăng cường độ
dốc (gradient boosting) cho phân vị 90, cùng mô hình đó nhưng bỏ số khách khác trên xe, và hồi quy phân vị tuyến tính.
Mỗi mô hình được so với một cách đệm đều có cùng độ dài thêm trung bình 30 giây, trên các ngày kiểm không dùng để học.
Một mô hình chỉ được coi là tốt hơn khi thỏa cả bốn điều kiện: tỉ lệ trễ quá một phút thấp hơn; độ dài thêm trung
bình không quá 1 giây so với đối chứng; phân vị 90 của độ dài thêm không quá 15 giây so với đối chứng; tỉ lệ tới sớm
quá một phút tăng không quá 1 điểm phần trăm.
% EV: G01; N4B-08 (MLS:26-35, 41-61, 63-80, 91-95); G02; G03

Không mô hình nào thỏa đủ bốn điều kiện. Với luật C có dự báo, mô hình không giảm thêm tỉ lệ trễ. Với điều phối không
trần, mức giảm nhiều nhất là 0,69 điểm phần trăm (145 khách), kèm lời hứa dài hơn mức cho phép. Với ghép tự do trên
tập tắc đường cố định, tỉ lệ trễ giảm từ 12,89\% xuống 11,65\% trên 807 khách, nhưng 7,1\% khách tới sớm quá một phút,
so với 0\% khi đệm đều. Bỏ số khách trên xe không đổi kết quả nào.
% EV: G05; G06; G07; G08; G09; G10 (câu được phép, HO:346-347); N4B-08 (ML:6,14,15,21: n và các điều kiện không thỏa); K40

% -----------------------------------------------------------------------------------------------------
\section{So với cách đệm lời hứa}
\label{muc:c4-dem}

Khi có dự báo, cộng thêm 30 giây vào mọi lời hứa của điều phối không trần cũng cho 2.328 lượt không trễ quá một
phút trên 3.672 yêu cầu, bằng luật C trần 30 giây. Ở ngưỡng hai phút, cách đệm cho 2.366 lượt, so với 2.338 của luật
C. Đổi lại, lời hứa dài thêm trung bình 37,2 giây thay vì 7,4 giây.
% EV: C10; N-16; N-18; K06

Đệm 30 giây ở đây là cộng đúng 30 giây vào lời hứa đầu theo dự báo của mọi khách, phát lại trên các lần chạy của
điều phối không trần: quyết định, tập khách và giờ trả thực tế giữ nguyên. Vì vậy nó không phải đối chứng ghép cặp
với luật C; hai bên chỉ chung 2.207 khách. Bảng~\ref{bang:c4-dem} xếp các cách hứa thành ba khối có độ dài thêm gần
bằng nhau.
% EV: C10 (lưu ý: phát lại, không phải đối chứng ghép cặp); ADDENDUM-1 phần đầu (quy tắc đệm D = D0 + k); C04; R25 §6

\begin{table}[htbp]
\centering
\caption{Các cách hứa trên tập kiểm tra chính khi có dự báo (34 kịch bản, 3.672 yêu cầu mỗi cấu hình; dự báo chỉ có
ở 21 kịch bản ngày thường). Số đếm gộp; trễ tính trong số khách được phục vụ, so với lời hứa của chính cách hứa đó;
hai cột ``không trễ quá'' tính trên 3.672 yêu cầu; độ dài thêm là trung bình gộp theo khách được phục vụ.}
\label{bang:c4-dem}
\small
\begin{tabular}{@{}p{4.4cm}rrrrr@{}}
\toprule
\textbf{Cách hứa} & \textbf{Phục vụ} & \textbf{Trễ quá 60 s} & \textbf{Không trễ quá 60 s} & \textbf{Không trễ quá 120 s}
& \textbf{Dài thêm (s)}\\
\midrule
\multicolumn{6}{@{}l}{\textit{Không đệm (dài thêm khoảng 7 giây)}}\\
U, dự báo & 2.434 & 137 & 2.297 & 2.347 & 7,18\\
C-30, dự báo & 2.338 & 10 & 2.328 & 2.338 & 7,37\\
C-120, dự báo & 2.387 & 65 & 2.322 & 2.379 & 7,27\\
% EV: N-16; N-03; C09; N-06; L1a-03
\midrule
\multicolumn{6}{@{}l}{\textit{Đệm 30 giây (dài thêm khoảng 37 giây)}}\\
U, dự báo, đệm 30 s & 2.434 & 106 & 2.328 & 2.366 & 37,18\\
C-30, dự báo, đệm 30 s & 2.338 & 2 & 2.336 & 2.338 & 37,37\\
% EV: C10; N-02; N-16
\midrule
\multicolumn{6}{@{}l}{\textit{Đệm 60 giây (dài thêm khoảng 67 giây)}}\\
U, dự báo, đệm 60 s & 2.434 & 87 & 2.347 & 2.386 & 67,18\\
% EV: C11; N-16
\bottomrule
\end{tabular}
\end{table}

Cách đọc bảng: khi đệm 30 giây trên điều phối không trần, 106 trong 2.434 khách được phục vụ tới muộn hơn lời hứa
đã đệm quá 60 giây, nên có 2.328 lượt không trễ quá 60 giây trên 3.672 yêu cầu.
% EV: C10; N-16

\textbf{Vai trò của từng thành phần.} Thứ nhất, ở độ dài thêm gần bằng nhau, luật C hơn ở ngưỡng một phút và kém ở
ngưỡng hai phút. Khi chưa đệm, luật C có 2.328 so với 2.297 lượt ở ngưỡng một phút và 2.338 so với 2.347 lượt ở ngưỡng
hai phút; khi cùng đệm 30 giây, các cặp số là 2.336 so với 2.328 và 2.338 so với 2.366. Cách đệm chỉ bằng luật C ở
ngưỡng một phút khi lời hứa dài thêm khoảng năm lần (37,18 so với 7,37 giây).
% EV: N-17; N-18 (37,1833 / 7,3671 = 5,05); N-16

Thứ hai, lời hứa dài thì ít trễ nhưng nhiều khách tới sớm (Bảng~\ref{bang:c4-haiphia}). Với luật C có dự báo, 4
trong 2.338 khách tới sớm quá một phút; với cách đệm 30 giây là 44 trong 2.434 khách. Số lượt tới trong khoảng
$\pm$60 giây quanh lời hứa là 2.324 so với 2.284 trên 3.672 yêu cầu.
% EV: C10; N-04; N4B-09 (ADD1:14-26,84-89)

\begin{table}[htbp]
\centering
\caption{Hai phía của độ lệch so với lời hứa khi có dự báo, trên tập kiểm tra chính (34 kịch bản). Các cột trễ và
tới sớm đếm trong số khách được phục vụ (2.434 với U, 2.338 với C-30, 2.387 với C-120); cột ``trong $\pm$60 s''
đếm số lượt tới cách lời hứa không quá 60 giây, tính trên 3.672 yêu cầu.}
\label{bang:c4-haiphia}
\small
\begin{tabular}{@{}p{4.4cm}rrrrr@{}}
\toprule
\textbf{Cách hứa} & \textbf{Tới sớm quá 60 s} & \textbf{Trong $\pm$60 s} & \textbf{Trễ quá 120 s} &
\textbf{Trễ quá 300 s} & \textbf{Trễ quá 600 s}\\
\midrule
U, dự báo & 8 & 2.289 & 87 & 17 & 2\\
C-30, dự báo & 4 & 2.324 & 0 & 0 & 0\\
C-120, dự báo & 2 & 2.320 & 8 & 0 & 0\\
U, dự báo, đệm 30 s & 44 & 2.284 & 68 & 11 & 1\\
C-30, dự báo, đệm 30 s & 44 & 2.292 & 0 & 0 & 0\\
U, dự báo, đệm 60 s & 1.033 & 1.314 & 48 & 8 & 1\\
\bottomrule
\end{tabular}
\end{table}
% EV: N-01; N-02; N-03; N-04; N-06; C11; N4B-09 (ADD1:14-16,19-20,24,84-89; A:7,9: late600)

Cách đọc bảng: với cách đệm 60 giây, 1.033 trong 2.434 khách tới sớm hơn lời hứa quá một phút, và chỉ 1.314 lượt tới
trong khoảng $\pm$60 giây quanh lời hứa.
% EV: C11; N4B-09 (ADD1:16,89)

Thứ ba, với cách đệm 30 giây, 68 khách vẫn trễ quá hai phút và 1 khách trễ quá mười phút; luật C có dự báo không có
khách nào trễ quá hai phút. Lợi thế 28 lượt của cách đệm ở ngưỡng hai phút đi cùng 96 lượt phục vụ nhiều hơn, không
phải ít khách trễ hơn.
% EV: N-01; N-19

Thứ tư, cách đệm không đổi quyết định điều phối, nên giữ nguyên 174 khách bị điều phối dời giờ trả, 12 khách bị dời
từ hai lần và một khách bị dời tổng cộng 590,1 giây. Luật C trần 30 giây giới hạn phần này của mỗi khách, nhiều nhất
29,2 giây (Bảng~\ref{bang:c4-doidieuphoi}), kể cả ở các kịch bản cuối tuần không có dự báo và ở chiều tuyết.
% EV: C20; N-23; C10 (lưu ý: đệm là phát lại trên quỹ đạo của U); N-28; N-26; ADDENDUM-1 mục 0.4

So với luật hạn chót 30 giây có dự báo, luật C có cùng 2.338 lượt phục vụ và cùng số lượt không trễ quá một và hai
phút (2.328 và 2.338). Khác biệt là luật hạn chót phải phục hồi ở 41,3\% lượt quyết định (trung bình theo kịch bản)
và phá lời hứa của chính nó ở 28 trong 34 kịch bản, còn luật C không lần nào (Bảng~\ref{bang:c4-nghiavu}).
% EV: N-16 (lưu ý: M⁺-30 + dự báo 2.338 / 10 / 2.328 / 2.338, ADD1:29); C13; N-21; X3-16

Trần cũng là một núm vặn. Trần 120 giây có dự báo cho 2.379 lượt không trễ quá hai phút, nhiều hơn cách đệm 30 giây,
với lời hứa dài thêm 7,27 giây, khoảng một phần năm, và phục vụ ít hơn điều phối không trần 47 lượt. Đây là một điểm
đo thêm trên đường đánh đổi, trên một tập khách khác, không phải mức trần tốt nhất.
% EV: N-06; N-18 (0,1956; 47); N-10; K36

\textbf{Khi nào chọn trần, khi nào chọn đệm.} Nếu hãng xe chỉ tính số lượt không trễ quá hai phút, và chấp nhận lời
hứa dài hơn cùng nhiều khách tới sớm hơn, đệm trên điều phối không trần là lựa chọn cạnh tranh và phục vụ thêm 96
lượt. Nếu cần giới hạn phần dời do ghép thêm cho từng khách, giữ được cả khi không có dự báo, cùng lời hứa ngắn, thì
luật C cho điều đó; đổi lại là 96 lượt phục vụ và 242 khách đi chung. Hai cách dùng chung được, như hàng luật C có dự
báo và đệm 30 giây. Khóa luận không có bằng chứng khách hay hãng xe coi trọng các điều này đến mức nào, nên đây là lựa
chọn của hãng xe, không phải kết luận của thực nghiệm.
% EV: N-19; C02; N-20; N-02; D20; I07; I10; K38; K51

% -----------------------------------------------------------------------------------------------------
\section{Khi dự báo không đủ: chiều tuyết 15/11/2018}
\label{muc:c4-tuyet}

Trên 601 khách được phục vụ của chiều tuyết 15/11/2018, dự báo lịch sử chỉ giảm số khách tới muộn hơn lời hứa đầu
quá một phút từ 41 xuống 34. Tính gộp theo khách, tỉ lệ giảm từ 6,82\% xuống 5,66\%; tính trung bình theo kịch bản,
từ 7,19\% xuống 5,94\%.
% EV: F13; K25; B20

Tập này gồm 8 kịch bản khung 14 và 15 giờ, chạy với luật C trần 30 giây; bộ chạy của tập chưa được kiểm lại. Theo sổ,
trong 3.445,9 giây trễ của 41 khách, 8,6 giây tính cho điều phối, 251,5 giây được hồ sơ lịch sử giải thích và 3.185,9
giây (92,5\%) chưa được giải thích (Hình~\ref{hinh:c4-so}).
% EV: B20; F13; N4B-05; N4B-06 (rider_nowcast_replay.py:6: 8 lần chạy C-30 không dự báo); phần A Bảng tập thử (bộ chạy chưa kiểm lại)

Bảng~\ref{bang:c4-tuyet} phát lại thêm vài cách hứa dùng thông tin cùng ngày. Các cách ``trộn'' lấy trung bình nhân
giữa mức tăng theo hồ sơ lịch sử và xu hướng tăng gần nhất của hệ số tắc đường trong ngày. Con số 19 dùng dữ liệu
buổi sáng không được giao cho khung này; một quy tắc chuyển chế độ có ngưỡng chọn sau khi xem dữ liệu cũng cho đúng
con số đó. Con số 29 vẫn lấy mức hiện tại đúng từ bảng thời gian đi.
% EV: F14; K20; K21; N4B-06 (factor_nowcast_screen.py:12-30; factor_nowcast_screen_v2.py:6-13; RP1:7-20; NEW:5)

\begin{table}[htbp]
\centering
\caption{Phát lại các cách hứa trên 8 kịch bản chiều tuyết 15/11/2018 (luật C trần 30 giây, 601 khách được phục
vụ). Cột ``khách trễ'' đếm gộp số khách tới muộn hơn lời hứa đầu quá 60 giây; cột cuối là độ dài thêm trung bình
theo kịch bản. Mọi cách hứa đều tính lại sau khi chạy, trên ngày đã được xem.}
\label{bang:c4-tuyet}
\small
\begin{tabular}{@{}p{3.6cm}p{7.0cm}rr@{}}
\toprule
\textbf{Cách hứa} & \textbf{Thông tin dùng lúc hứa} & \textbf{Khách trễ} & \textbf{Dài thêm (s)}\\
\midrule
Không dự báo & Thời gian đi hiện tại & 41 & 0\\
Dự báo lịch sử & Hồ sơ tắc đường của các ngày nguồn & 34 & 2,0\\
Trộn trong khung & Thêm hệ số cùng ngày đã có trong khung giờ & 21 & 8,5\\
Trộn có buổi sáng & Thêm hệ số cùng ngày từ 6 giờ, không được giao cho khung này & 19 & 9,5\\
Trộn trễ một khung & Như hàng trên, xu hướng chỉ dùng các khung đã qua & 29 & 9,0\\
Đệm cùng giá & Tỉ lệ đệm lấy theo hàng ``trộn trong khung'' trên chính ngày tuyết & 27 & 8,3\\
Hệ số thật & Hệ số sau lúc hứa, không có ở thời điểm hứa & 0 & 18,8\\
\bottomrule
\end{tabular}
\end{table}
% EV: F14 (RP1:7-20; NEW:5); N4B-06 (định nghĩa các cách hứa; cột dài thêm là trung bình theo kịch bản, RP1:6, rider_nowcast_replay.py:96-104)

Cách đọc bảng: khi trộn hồ sơ lịch sử với xu hướng trong khung giờ, 21 trong 601 khách tới muộn hơn lời hứa quá một
phút, với lời hứa dài thêm trung bình 8,5 giây.
% EV: F14; N4B-06

Theo sổ, mô hình có dữ liệu buổi sáng giải thích thêm 886,9 giây (25,7\%), còn 2.299,0 giây (66,7\%) chưa được giải
thích bằng các cách đã thử. Đây là một sự kiện đã được xem trước. Giai đoạn 9/2018--12/2019 chỉ có 5 đợt tuyết từ
25,4 mm trở lên, kể cả đợt này, nên một phép kiểm dự báo cho ngày có sự kiện cần nhiều sự kiện độc lập hơn.
% EV: F15; K19; I11; B26

% -----------------------------------------------------------------------------------------------------
\section{Kiểm tra tính hợp lệ của kết quả}
\label{muc:c4-hople}

Trên 34 kịch bản của tập kiểm tra chính, không khách nào vượt trần: tổng dời do điều phối lớn nhất của một khách là
29,2 giây với trần 30 giây và 118,9 giây với trần 120 giây. Không có lượt quyết định nào phải phục hồi, 0 trong 20.101
và 0 trong 20.435 lượt.
% EV: C20; N-07; X3-16; N-21

Giới hạn tính toán không cắt kết quả trên tập này. Ở bốn cấu hình được kiểm (điều phối không trần, luật C trần 30 và
120 giây, luật hạn chót 30 giây), giới hạn 100 phương án mỗi xe chưa lần nào bỏ phương án hợp lệ. Giới hạn khối lượng
khám phá chưa lần nào để sót nhánh chưa xét, và giới hạn 2 yêu cầu mới mỗi lượt chỉ một lần để một yêu cầu sang lượt
sau trong 34 lần chạy của mỗi cấu hình. Chênh lệch giữa các cấu hình vì vậy phản ánh tác động của cả gói cam kết,
không phải của giới hạn tính toán. Bộ sinh vẫn không vét cạn: nó chỉ chèn yêu cầu mới vào kế hoạch đang có
(mục~\ref{muc:gioihan}).
% EV: N-25 (ADD1:119-122; A:7-9); N4B-11 (GEN:54-95, 298-345; CandidateModels.cs:100-118; R15-audit.py:76-79); K34; X3-03; X3-32; A24

Phần đầu của tiêu chí 3 không phân biệt phương án trên tập này. Ở mọi lượt có ghi tiêu chí này (3.503 lượt khi không
trần, 3.506 với trần 30 giây, 3.504 với trần 120 giây), mức dùng trần lớn nhất đều bằng 100\%, vì đổi xe, đổi điểm đón
và đổi điểm trả có trần 0. Tiêu chí 3 vì vậy chỉ tác dụng qua tổng phần dời theo từng chiều.
% EV: A15; N-25; A07; X3-07

Tệp chính sách cam kết của điều phối không trần và của luật C trần 30, 120 giây trên tập kiểm tra chính giống từng
byte với tệp cùng tên của tập tắc đường cố định. Hai tập khác nhau ở thiết lập bộ chạy: luật phục hồi, cách ghi lần
đón muộn, phiên bản chính sách và mức ghi chép của bộ giải.
% EV: A08; NEW:11-13; N4B-12

Các số của tập kiểm tra chính được ít nhất hai bộ đọc viết độc lập tính lại từ bản ghi thô và khớp trong phạm vi đã
kiểm, gồm cả từng khách ở 7.159 khách được phục vụ của điều phối không trần và luật C trần 30, 120 giây. Mục~\ref{muc:kiemtra}
nêu thêm phép kiểm 382.006 bản ghi sửa và 13 hạng mục được tính lại. Khi chạy lại ngày 12/10/2026 trên đúng mã nguồn
đã dựng bộ chạy, với kiểu xuống dòng LF, 1.024 trong 1.025 kiểm thử tự động đạt. Kiểm thử còn lại đòi cây làm việc có
tệp chưa được ghi nhận, và đạt khi chạy riêng với một tệp như vậy.
% EV: N4B-14 (ADD1:66-78); K41; A26; C19; N4B-13 (TESTRUN-45f8026-2026-10-12.md; log dòng 43-105; log dirtycheck)

Các phép kiểm trên cho biết con số khớp bản ghi và cài đặt giữ đúng trần trên các lần chạy đã kiểm. Chúng không cho
biết mô phỏng giống thực tế đến đâu (mục~\ref{muc:c4-thaoluan}).
% EV: I01; K41

% -----------------------------------------------------------------------------------------------------
\section{Thảo luận và hạn chế}
\label{muc:c4-thaoluan}

Với một hãng xe, các kết quả của chương là một bảng các điểm đánh đổi, không phải một khuyến nghị. So với điều phối
không trần, luật C trần 30 giây đổi 96 lượt phục vụ và 242 khách đi chung lấy ba điều: phần dời do điều phối của mỗi
khách không quá 30 giây, số lượt trễ thấp hơn ở cả bốn ngưỡng đã đo, và không lượt quyết định nào phải phục hồi. Trần
120 giây là một điểm ở giữa; cách đệm phục vụ nhiều hơn với lời hứa dài hơn. Khóa luận không đo giá trị của các điều
này với khách hay hãng xe.
% EV: C02; N-20; C20; N-24; X3-16; N-06; C10; D20; I10; K38

Một phép tính giả định cho thấy cách đọc bảng này. Giả sử mỗi lượt trễ quá một phút làm hãng xe tốn một khoản $k$ và
mỗi chuyến mang lại biên lợi nhuận $m$. Với 96 lượt phục vụ ít hơn và 127 lượt trễ ít hơn, luật C chỉ có lợi hơn khi
$k > 0{,}756\,m$ (tức $96/127$), nếu mọi yếu tố khác bằng 0. Hai nhóm khách khác nhau, nên đây không phải phép đổi một
lượt lấy một lượt; khóa luận cũng không có giá trị thật nào của $k$ hay $m$. Một giả thuyết chưa kiểm: ở tải cao hơn,
mỗi lần ghép thêm chạm tới nhiều khách đang trên xe hơn, nên trần có thể chặn nhiều phương án và mất nhiều lượt phục
vụ hơn.
% EV: C29 (R15:231-235; 96/127 = 0,756); K35; C28; I10; câu cuối là giả thuyết, không có dòng sổ (RD LR-34); B12

Các hạn chế dưới đây chia theo bốn nhóm; mỗi ý nêu điều chưa làm và kết luận chịu ảnh hưởng.

\textbf{Mô hình mô phỏng.}
\begin{itemize}[leftmargin=1.2cm,itemsep=2pt,topsep=3pt]
\item Mọi kết quả đến từ một bộ mô phỏng, FleetPy, nên không cho biết hiệu quả khi vận hành thật.
\item Tắc đường là một hệ số chung cho cả thành phố theo khung 15 phút, mức nền lấy lúc 00:00. Thế giới vì vậy gần
như tất định, và mức trễ còn lại sau dự báo không chuyển sang mạng có tắc đường riêng trên từng đoạn đường.
\item Thời gian dừng đón, trả bằng 0 giây (kiểm trên 3 lần chạy mẫu), trong khi nghiên cứu điển hình của FleetPy dùng
30 giây mỗi khách~\cite{engelhardt2026}; đồ thị là bao đóng 96 nút, không phải mạng đường phố. Độ dài chuyến và độ
lệch giờ trả chưa tính các yếu tố này.
\item Không có phản ứng của khách (hủy chuyến, từ chối lời mời) và không cho chuyển khách đã nhận sang xe khác. Số
lượt phục vụ vì vậy không tính khách bỏ đi, và chưa biết cái giá của trần đổi thế nào khi được gán lại xe.
\end{itemize}
% EV: I01; B01; B06; B07; I02; F10; B05 (SERV:2-4; FLEET:423-424); J18; B03; K48; I07; X3-39

\textbf{Dữ liệu và thiết kế thực nghiệm.}
\begin{itemize}[leftmargin=1.2cm,itemsep=2pt,topsep=3pt]
\item Tập kiểm tra chính có 34 kịch bản phụ thuộc nhau thuộc 5 ngày, mỗi kịch bản chạy một lần. Các phép đếm theo
kịch bản là mô tả, không phải kiểm định, và không bảo đảm cho ngày khác.
\item Ngưỡng phân tích, cách chia tập và mức 120 giây được chọn sau khi xem dữ liệu phát triển; chiều tuyết là một sự
kiện đã xem. Mọi kết quả là thăm dò.
\item Tải thấp: trung bình 0,535 khách trên xe kể cả lúc xe rỗi với điều phối không trần, so với khoảng 1,33 trong
nghiên cứu điển hình của FleetPy (180 xe, 21.400 yêu cầu). Các chuyến được chọn cũng ngắn hơn tập đủ điều kiện (trung
vị đi thẳng 165 so với 196 giây). Kết quả không suy ra cho tải cao hay chuyến dài hơn.
\end{itemize}
% EV: B09; I05; B11; I06; B10 (lưu ý); I11; B12 (LOAD:2-3; FLEET:433-434,447-449); J18; B13 (SEL6:2-4)

\textbf{Tích hợp.}
\begin{itemize}[leftmargin=1.2cm,itemsep=2pt,topsep=3pt]
\item Lời hứa theo dự báo được tính sau khi chạy, còn phép kiểm khả thi vẫn dùng giờ không dự báo; cách đưa dự báo
vào bộ điều phối đã thử làm 3 trong 32 lần chạy dừng. Muốn triển khai phải sửa phép kiểm giờ đón trước.
\item Bộ chạy dựng lại từ mã nguồn có thể khác bản lưu trữ, và bộ chạy của chiều tuyết chưa được kiểm lại. Lần chạy
mới phải cho đúng kết quả cũ trên vài kịch bản trước khi đem so với các số ở đây.
\end{itemize}
% EV: F07; F12; I13; R18 B24 (qua REG A28); phần A Bảng tập thử (chiều tuyết: chưa kiểm lại)

\textbf{Phạm vi kết luận.}
\begin{itemize}[leftmargin=1.2cm,itemsep=2pt,topsep=3pt]
\item Bộ chọn ưu tiên chi phí chưa được chạy trên quỹ đạo (mục~\ref{muc:c4-vaitro}), và luật mỗi lần một trần chưa
được cài đặt (mục~\ref{muc:c4-congdon}). Chưa tách được phần của trần với phần của bộ chọn, và chưa biết vế cộng dồn
có đổi kết quả hay không.
\item Luật hạn chót chỉ là phần hạn chót, ở 30--60 giây thay vì 900 giây của Schulz và Pfeiffer~\cite{schulz2026};
hạn chót nới kiểu Hörl và Zwick~\cite{horl2022} chưa được chạy. Các phép so không nói về phương pháp của các tác giả
đó.
\item Mức 30 giây không phải ngưỡng cảm nhận hay chuẩn ngành, và không có dữ liệu về mức khách sẵn lòng trả cho lời
hứa ổn định hơn. Vì vậy không mức trần nào được chọn là tốt nhất.
\end{itemize}
% EV: H02; H12; K18; E09; K08; J10; J11; K27; D20; K52; K36
